(1) 令 $K=\sigma A$，辐射从高温物体1向低温物体2净流动，故 $C_1\dot T_1=-K(T_1^4-T_2^4)$、$C_2\dot T_2=+K(T_1^4-T_2^4)$。将方程相加可见 $C_1T_1+C_2T_2$ 不随时间变化，因未向环境流失能量；这里没有把热容错误视为每秒的热容流率。
(2) 长时间热流趋零，因为各温度均为正，$T_1^4=T_2^4$ 意味着两温度相等为 $T_f$。结合总能量守恒，$T_f=(C_1T_{1i}+C_2T_{2i})/(C_1+C_2)$，它严格位于两个初温之间。该结果由总能量决定，与交换面积只影响速度的事实相符。
(3) 每个恒热容物体的熵变化为 $\Delta S_j=C_j\ln(T_f/T_{ji})$，因此 $\Delta S_{\rm total}=C_1\ln(T_f/T_{1i})+C_2\ln(T_f/T_{2i})$。沿时间求导更直接：$\dot S=K(T_1^4-T_2^4)(1/T_2-1/T_1)\ge0$，因两个括号同号。由此从初始至最终积分证明总熵变化非负，严格不等号对应初温不同。
(4) 令 $\Delta=T_1-T_2$，在两温度均接近 $T_f$ 时有 $T_1^4-T_2^4=4T_f^3\Delta+O(\Delta^2)$。两动力学式相减，$\dot\Delta=-4\sigma A T_f^3(1/C_1+1/C_2)\Delta+O(\Delta^2)$。因此温差按 $e^{-t/\tau}$ 衰减，$\tau=[4\sigma A T_f^3(1/C_1+1/C_2)]^{-1}$，体现高温时辐射换热加快。
